\section{Conclusion}
\label{sec:conclusion}

The lowest test sMAPE depends on the grain.
Pooling the four ranks into one winner would contradict Table~\ref{tab:top3}.

On the monthly panel, E2 has the lowest test sMAPE, about 10.9.
Boosting (about 13.3) and the lag-12 naive (about 13.5) are close behind.
E1 does not lead that grain.
E2 switches on absolute disagreement.
E1 switches on an energy threshold.
The two rules are not the same model.

On calendar-daily A, TimesFM has the lowest test sMAPE, about 59.8.
Most other models have negative skill against the lag-7 naive.
Zeros from weekends and empty slips are inside that score.
sMAPE on days with positive collection is a side description.
It is not Ranking A.

On business-daily B, TimesFM again has the lowest test sMAPE, about 52.4.
Boosting and least squares follow.
Many positive skill scores on B are skill against lag 5.
Lag 5 is a weak reference on this series.
National holidays were kept, and no model received a holiday indicator.

On the weekly series, TimesFM has the lowest test sMAPE, about 17.2, and skill about 0.41 against the lag-52 naive.
Ridge and least squares are next.

The seasonal naive remains the reference forecast on each grain.
E1 is the rule that can abstain by energy.
Reading E1 requires coverage next to sMAPE, and this extract does not put coverage in the rank tables.
E2 is the disagreement switch.
TimesFM has the lowest test sMAPE on A, on B, and on the weekly series.
It does not have the lowest test sMAPE on the monthly panel.

Later work on the same extract should publish E1 coverage with the sMAPE tables, add a holiday indicator on grain B, and keep each grain's own seasonal naive as the skill reference.
